\scanpage{348}
$\unsure{\frac12|k}$-Algèbre différentielle $C^*$ sur Topos $\mathcal X$.

Anneau de base $k$ (commutatif)
résolution $C^*$ de $k$

\margin{\unsure{ou, il suffit} \unsure{qu'il y ait une} \unsure{augmentation|graduation}}

\diagram{\begin{tikzcd}
\cdot\,\cdot \arrow[r,dash] & C^0 \arrow[r,dash] & C^1 \arrow[r] & C^2 \arrow[r] & -\\
& k \arrow[u]
\end{tikzcd}}

(la section unité de $k$ correspond à une section $\unsure{T|\overline1}$ de $C^0$).

\deleted{\unsure{bilin.|bilin}}

Accouplements (notés $x,y\longmapsto x\circ y$)
\[
\left\{\begin{array}{lll}
Z^i\times C^j\longrightarrow C^{i+j}&(a_{ij})\\
C^i\times Z^j\longrightarrow C^{i+j}&(b_{ij})&i,j\unsure{\leq|\geq}0
\end{array}\right.
\]
satisfaisant aux conditions

$1^\circ)$ $a_{ij}$ et $b_{ij}$ coïncident sur $Z^i\times Z^j$ et
y prennent valeurs dans $Z^{i+j}$. (Donc la notation
$x\circ y$ non ambiguë, elle a un \deleted{\unsure{donc}} sens si $dx$ \uline{ou} $dy$ est \unsure{nulle|nul}.)

$2^\circ)$ Associativité $(x\circ y)\circ z=x\circ(y\circ z)$ \unsure{si deux}
des trois sections locales $x,y,z$ (de $C^i,C^j,C^h$)
sont des cocycles; $1\circ x=x\circ1=x$.

\[
3^\circ)\quad\begin{array}{l}
x\circ y=y\circ x\;(-1)^{\deg y\deg x}\\
\unsure{\rbrace|)}\quad x\circ x=0\quad\text{si $x$ section locale de $Z^i$, $i$ \uline{impair}}
\end{array}
\]

\deleted{si $\unsure{x^i|u^i}$ \unsure{section locale de} $C^i$, $\unsure{z^j|x^j}$ \unsure{section locale de} $\unsure{Z^j}$}

$4^\circ)$ Pour tout $i,j\geq1$, existe un $k$-bilin.
\diagram{\begin{tikzcd}
\inserted{\unsure{Z^i|Z^{i-1}}} & \\
\deleted{C^{i-1}\times Z^j}\arrow[r,"\unsure{h_{ij}|b_{ij}}"] & C^{i+j-1}
\end{tikzcd}}
rendant le carré commutatif :
\figure{\begin{tikzpicture}
\node (q) {\begin{tikzcd}
Z^{\unsure{i-1|i}}\times Z^j\arrow[r] & Z^{i+j}\\
C^{i-1}\times Z^j\arrow[u,"d\times\mathrm{id}"]\arrow[r,"\unsure{h_{ij}|b_{ij}}"] & C^{i+j-1}\arrow[u,"d"']
\end{tikzcd}};
\draw[thin] ([xshift=-2pt]q.west) -- ([xshift=-2pt,yshift=-2pt]q.south west) -- ([xshift=3pt,yshift=-2pt]q.south east) -- ([xshift=3pt,yshift=8pt]q.north east);
\end{tikzpicture}}
\margin{$\unsure{d(x^i\circ z^j)=dx^i\circ z^j}$}

\deleted{(donc \illegible{2} $3^\circ$, $\unsure{z^j\circ dx^i|z^j\circ dx^{i-1}}$
\deleted{\illegible{1}}
\[
\unsure{h_{ij}|b_{ij}}(dx^{i-1})\unsure{\wedge|\circ}z^j=d(\unsure{h_{ij}|b_{ij}}(x^{i-1},z^j)
\]
\illegible{3} \unsure{si} $x^{\unsure{i|i-1}}$ \unsure{section locale} \illegible{3})}

\deleted{(\uline{NB} l'\unsure{indétermination|interprétation} \unsure{des|de} $\unsure{h_{ij}|b_{ij}}$ \unsure{est dans}
\[
\operatorname{Hom}_k(C^{i-1}\otimes_k Z^j,Z^{i+j-1});
\]
\unsure{elle n'est pas} \unsure{voulue|variable}; on peut \unsure{exiger pour les}
\illegible{2} \unsure{des} \uline{\unsure{mieux|choix}}, et pour des \illegible{3}\unsure{...|....})}

\mspage{349}
